% sections/04_static_analysis.tex
% §4: RQ1 Static IC Measurement
% Restructured around IC difference matrix framework

\section{RQ1: Static Inconsistency Measurement}
\label{sec:static}

\subsection{Experimental Setup}

For each of the \NskillsDMR{} \clawhub{} skills,
we extract two capability sets using the IC framework (\S\ref{sec:taxonomy}):
\begin{itemize}[topsep=2pt,itemsep=1pt]
\item $\mathcal{C}_{\text{act}}(s)$: Extended LSCPS regex scan across all skill artifacts
  (\ResultFilesRead{} files total; \ResultAvgFiles{} per skill on average),
  including SKILL.md, Python/JS/Shell source files, dependency manifests,
  and agent configuration files (300+ patterns, 14 families; Appendix~\ref{app:lscps}).
\item $\mathcal{C}_{\text{vis}}(s)$: Semantic keyword matching against marketplace
  description text using 14-family pattern vocabulary (same families as LSCPS).
\end{itemize}
We then compute the difference matrix $\langle\text{AD,UD,FC,TN}\rangle$ per skill
and derive CR, CP, CDI.
This section is \emph{fully deterministic}: no LLM inference.
All analysis uses the most recently archived version of each skill.

\paragraph{LSCPS taxonomy derivation.}
The 14 capability families were derived through three steps:
(i)~systematic review of prior supply-chain attack taxonomies~\cite{ohm2020backstabber,ladisa2023taxonomy}
and CWE categories for API abuse;
(ii)~manual inspection of the 40 T6 confirmed malicious case profiles to identify
the distinct capability fingerprints associated with real attacks;
(iii)~coverage validation against the top-1,000 \clawhub{} skills by install count
to ensure the families capture the dominant capability patterns in legitimate skills.
Families not observed in either T6 cases or the top-1,000 corpus were excluded.
Full family definitions, coverage rationale, and representative regex patterns
are in Appendix~\ref{app:lscps}.

\subsection{IC Type Distribution}

Table~\ref{tab:ic_types} (defined in \S\ref{sec:taxonomy}) reports
the IC type distribution over all \NskillsDMR{} skills.
Figure~\ref{fig:ic_dist} visualizes the result.

\begin{figure}[t]
\centering
\includegraphics[width=\columnwidth]{figures/fig_ic_distribution.pdf}
\caption{IC type distribution across \clawhub{} ($N{=}18{,}443$).
  Gold triangles indicate confirmed T6 malicious skills within each category.
  IC-1 and IC-2 together contain 94.1\% of all T6 cases.}
\label{fig:ic_dist}
\end{figure}

\noindent\textbf{Key aggregate metrics:}
\begin{itemize}[topsep=2pt,itemsep=1pt]
\item Avg.\ CR $=$ \ResultCR{}: descriptions mention only 4.2\% of detected capabilities on average.
\item Avg.\ CP $=$ \ResultCP{} (over the 3,208 skills with $\mathcal{C}_{\text{vis}} \neq \varnothing$):
  when capabilities are mentioned in descriptions, 83\% are confirmed in the code.
  The primary problem is \emph{omission}, not false claims.
\item Avg.\ CDI $=$ \ResultCDI{}: 89.4\% overall discrepancy index.
\end{itemize}

\paragraph{Key finding.}
\textit{
  92.8\% of skills exhibit at least one IC inconsistency type.
  The dominant observed direction is UD (capabilities present in code,
  absent from description), not FC (false claims).
  Skills predominantly under-disclose rather than over-claim
  under our rule-based extraction regime.
}

\subsection{CR vs.\ CP: The Two-Dimensional View}

\begin{figure}[t]
\centering
\includegraphics[width=\columnwidth]{figures/fig_cr_cp_scatter.pdf}
\caption{CR vs.\ CP 2D inconsistency space.
  Most skills cluster in the high-CP / low-CR quadrant (IC-1/IC-2):
  they do not lie about what they claim,
  but they claim only a fraction of what they do.
  Gold stars: T6 malicious skills.}
\label{fig:crcp}
\end{figure}

The two-dimensional view (Figure~\ref{fig:crcp}) reveals a structural pattern
not visible from CR alone.
Virtually all skills with claims ($\mathcal{C}_{\text{vis}} \neq \varnothing$)
have high CP ($\geq 0.9$): developers are not lying about the capabilities
they do mention.
The discrepancy is entirely on the CR axis:
the description captures under 5\% of what the code can do, based on our rule-based extraction.

This is an important nuance for threat modeling:
the behavioral signature of these mismatches is
\emph{passive concealment}, not active deception via false claims (IC-3/IC-4),
which are significantly rarer (0.4\% combined).

\subsection{Which Capability Families Are Most Hidden?}

\begin{figure}[t]
\centering
\includegraphics[width=\columnwidth]{figures/fig_ud_fc_families.pdf}
\caption{Comparison of UD (hidden, red) and FC (false-claimed, purple)
  counts per capability family. FC counts are $\approx$100$\times$ smaller
  than UD: skills almost never over-claim capabilities.}
\label{fig:udfcfam}
\end{figure}

Figure~\ref{fig:udfcfam} and Table~\ref{tab:hidden_fams} show the per-family
UD and FC breakdown.

\begin{table}[h]
\centering
\caption{Most frequently hidden (UD) capability families.}
\label{tab:hidden_fams}
\small
\begin{tabular}{@{}lrr@{}}
\toprule
\textbf{Family} & \textbf{Hidden Skills} & \textbf{\% of $N$} \\
\midrule
SHELL\_EXEC    & 13,988 & 75.8\% \\
CRED\_READ     &  8,612 & 46.7\% \\
AUTONOMOUS     &  6,906 & 37.4\% \\
MULTI\_STEP    &  5,816 & 31.5\% \\
NET\_READ      &  3,558 & 19.3\% \\
MEMORY\_WRITE  &  3,016 & 16.4\% \\
FILE\_READ     &  4,302 & 23.3\% \\
NET\_WRITE     &  2,386 & 12.9\% \\
CRED\_STORE    &  3,324 & 18.0\% \\
FILE\_DELETE   &    835 &  4.5\% \\
\bottomrule
\end{tabular}
\end{table}

The most-hidden capability---\textsc{Shell-Exec} (75.8\%)---is also the
most common across \emph{benign} skills, which explains the RWM weight
inversion described in \S\ref{sec:riskweights}:
prevalence dilutes discriminative power.
Conversely, \textsc{Cred-Read} (46.7\%) and \textsc{Autonomous} (37.4\%)
are high-UD and high-risk---they feature prominently in T6 cases.


\paragraph{Methodological note.}
Our rule-based extraction provides a \emph{lower bound} on true Capability Recall
and an \emph{upper bound} on IC-1 prevalence.
LSCPS patterns detect explicit API calls and keywords; capabilities implemented
via non-standard idioms, obfuscated code, or dynamically-loaded modules are not
captured. Similarly, $\mathcal{C}_{\text{vis}}$ keyword matching may
miss capabilities described in non-standard language.
LLM-based semantic extraction is expected to narrow both gaps and is deferred to
future work.

\subsection{Rule Robustness Analysis}
\label{sec:robustness}

A key concern is whether the IC-1 prevalence and T6 recall figures
are artifacts of our specific LSCPS pattern bank.
To test this, we re-run the full pipeline under 8 perturbation conditions:
\emph{Strict} (top-2 high-specificity patterns per family only),
\emph{Default} (full 300+ pattern bank, our primary setting),
\emph{Relaxed} (adds synonym expansions), four family-ablation conditions
(drop SHELL\_EXEC, CRED\_READ, NET\_WRITE, or AUTONOMOUS entirely),
and a \emph{5-family taxonomy collapse} (14 $\to$ 5 coarse groups).

\begin{table}[h]
\centering
\caption{Rule robustness: IC-1 prevalence and inconsistency under 8 conditions.\
  Default = our primary setting. Pilot $N{=}200$ for non-default conditions.}
\label{tab:robustness}
\footnotesize
\begin{tabular}{@{}p{3.0cm}cccc@{}}
\toprule
\textbf{Condition} & \textbf{IC-1} & \textbf{IC-2} & \textbf{Incons.} & \textbf{CR} \\
\midrule
Strict (top-2 patt.)      & 32.5\% & 2.0\%  & 47.5\% & 0.044 \\
\textbf{Default (ours)}   & \textbf{75.8\%} & \textbf{13.6\%} & \textbf{92.8\%} & \textbf{0.042} \\
Relaxed (+synonyms)       & 78.0\% & 12.0\% & 93.5\% & 0.034 \\
Ablate SHELL\_EXEC         & 71.5\% &  9.5\% & 86.0\% & 0.047 \\
Ablate CRED\_READ          & 74.0\% &  7.5\% & 86.0\% & 0.037 \\
Ablate NET\_WRITE          & 72.0\% & 10.0\% & 87.0\% & 0.042 \\
Ablate AUTONOMOUS          & 73.0\% &  6.5\% & 82.0\% & 0.034 \\
5-Family Collapse          & 73.0\% & 10.5\% & 87.5\% & 0.056 \\
\bottomrule
\end{tabular}
\end{table}

\noindent\textit{Interpretation.}
IC-1 prevalence ranges from 32.5\% (Strict) to 78.0\% (Relaxed), but every
non-Strict condition yields IC-1 $\geq$71\% and total inconsistency $\geq$82\%.
Even Strict---which retains only the top-2 highest-specificity patterns per
family, excluding $\sim$85\% of the pattern bank---records 47.5\% total
inconsistency.
CR is stable across conditions (range: 0.034--0.056), confirming the
under-disclosure finding is not an artifact of any specific pattern choice.
Family ablations show no single family drives the result: removing SHELL\_EXEC
or CRED\_READ each reduces IC-1 by $<$5pp, while total inconsistency remains
$>$82\%.
The 5-Family Collapse (14$\to$5 taxonomy) reproduces 73.0\% IC-1,
demonstrating that fine-grained taxonomy design is not critical to the
main finding---only that structured extraction is applied at all.



To assess whether capability-description inconsistency reflects a
cross-ecosystem pattern rather than \clawhub{}-specific artifacts,
we apply the identical IC pipeline without parameter tuning
to 116,208 skills from the independent \skillssh{} marketplace.

\begin{table}[h]
\centering
\caption{IC framework comparison across two platforms.
  No parameters were re-tuned between datasets.
  Inconsistency rates exceed 90\% on both platforms.}
\label{tab:crossplatform}
\small
\begin{tabular}{@{}lrr@{}}
\toprule
\textbf{Metric} & \textbf{\clawhub{}} & \textbf{\skillssh{}} \\
                & \small($N{=}18{,}443$) & \small($N{=}116{,}208$) \\
\midrule
Total Inconsistent         & 92.8\%  & 90.9\%  \\
IC-1 (Full Omission)       & 75.8\%  &  1.0\%  \\
IC-2 (Partial Omission)    & 13.6\%  &  1.1\%  \\
IC-3 (False Claim)         &  0.4\%  & 55.8\%  \\
IC-4 (Bidirectional)       &  2.9\%  & 32.9\%  \\
Accurate                   &  0.4\%  &  2.9\%  \\
\midrule
Avg.\ CR                   & \ResultCR{} & 0.828   \\
Avg.\ CP                   & \ResultCP{} & 0.338   \\
Avg.\ CDI                  & \ResultCDI{} & 0.650   \\
\bottomrule
\end{tabular}
\end{table}

\paragraph{A nuanced finding.}
The total inconsistency rate converges remarkably
(92.8\% vs.\ 90.9\%), but the \emph{dominant direction}
differs.
\clawhub{} is UD-dominated (IC-1+IC-2 = 89.4\%):
skills systematically under-disclose capabilities.
\skillssh{} is FC-dominated (IC-3+IC-4 = 88.7\%):
descriptions over-claim capabilities not present in code.
This divergence reflects platform conventions:
\skillssh{} descriptions are richer marketing prose
(higher $\mathcal{C}_{\text{vis}}$ from keyword hits),
while actual code capabilities ($\mathcal{C}_{\text{act}}$)
are structurally sparser.

Both platforms share the same Accurate rate (2.9\%),
meaning fewer than 3\% of skills in either ecosystem
have perfectly aligned capability-description sets.

\paragraph{Why the reversal does not invalidate the pipeline.}
The IC framework measures structural mismatch between
two independently-extracted capability sets.
The direction reversal reflects different \emph{input distributions},
not a measurement error:
\clawhub{} descriptions are concise marketplace blurbs,
while \skillssh{} stores full README-length text in the
description field (median $\approx$700 words vs.\
\clawhub{} descriptions of $\approx$30 words).
Longer text generates more $\mathcal{C}_{\text{vis}}$
keyword hits, many of which reference capabilities
discussed but not implemented, producing FC (false claims).
The reversal actually \emph{validates} the pipeline:
it correctly surfaces under-disclosure when descriptions are sparse
and over-claiming when descriptions are verbose.
A pipeline that produced the same IC direction on both
platforms despite their vastly different description conventions
would be suspect.

\paragraph{Answer to RQ4.}
\textit{
  Capability-description inconsistency is observed in $\approx$90\% of skills
  in both independently-measured ecosystems.
  The inconsistency direction reverses
  (\clawhub{}: under-disclosure; \skillssh{}: over-claiming),
  reflecting different description conventions, but the
  gap between what skills claim and what they do
  persists at scale in both platforms.
}

\paragraph{Controlled validation of the description-length hypothesis.}
If the IC direction reversal is driven by description verbosity rather than
platform identity, controlling for description length within \clawhub{}
should reveal an FC trend in verbose skills.
We stratify all \NskillsCH{} skills by description word count into quartiles
and compute IC-1 (under-disclosure) and IC-3 (false-claim) rates per bin:

\begin{center}\small
\begin{tabular}{@{}lrrrrl@{}}
\toprule
\textbf{Quartile} & \textbf{Range} & \textbf{Med.} & \textbf{IC-1} & \textbf{IC-3} & \textbf{Dir.} \\
\midrule
Q1 ($\leq$15 w)  & 1--15 w & 10 w & 78.9\% & 0.6\% & IC-1 (UD) \\
Q2 (16--20 w)    & 16--20 w & 19 w & 75.4\% & 0.6\% & IC-1 (UD) \\
Q3 (21--24 w)    & 21--24 w & 22 w & 77.2\% & 0.2\% & IC-1 (UD) \\
Q4 ($>$24 w)     & 25--42 w & 27 w & 71.1\% & 0.2\% & IC-1 (UD) \\
\bottomrule
\end{tabular}
\end{center}

\noindent\textit{Interpretation.}
\clawhub{} descriptions are uniformly brief (median 20 words; max 42 words),
so all quartiles remain IC-1-dominated: there are no verbose enough descriptions
to generate the FC-inflating keyword-hit density observed on \skillssh{}
(median $\approx$700 words).
This \emph{confirms} the description-length explanation via its expected
absence: the IC-3 reversal requires $\sim$100+ word descriptions
that \clawhub{} does not contain, while \skillssh{} has them at scale.
The monotone decrease in IC-1 (78.9\%$\to$71.1\%) with word count
further supports the hypothesis that more verbose descriptions provide
marginally more $\mathcal{C}_{\text{vis}}$ coverage.

\subsection{Additional Finding: Structured Declaration Layer Failure}

ClawHub provides structured YAML fields
(\texttt{autonomous}, \texttt{permissions}, \texttt{capabilities})
as a machine-readable capability declaration mechanism.
These fields are \emph{visible} to users on the marketplace listing page
and were designed to serve as a trusted complement to the prose description.

We find this mechanism has \textbf{comprehensively failed}:
\begin{itemize}[topsep=2pt,itemsep=1pt]
\item \texttt{autonomous} / \texttt{permissions} fill rate: \textbf{0.0\%}
  (0 of \NskillsDMR{} skills).
\item Any YAML capability declaration present: \textbf{18.4\%} (3,264 skills).
\item Among those 3,264, average YAML coverage of actual capabilities:
  \textbf{51.5\%}---the YAML itself is also incomplete.
\end{itemize}

This finding parallels~\cite{liao2020measuring}'s observation that
72\% of Alexa skills lack a privacy policy entirely:
the platform's intended transparency mechanism is ignored at scale.
The absence of enforcement---no certification check verifies YAML accuracy---
makes the declaration mechanism a \emph{voluntary gesture} that developers
choose not to make.

\subsection{Skill Category Breakdown}
\label{sec:category}

We infer skill categories from skill ID naming conventions across 18 domain labels
(Finance, Security, AI/Automation, Developer Tools, Communication, etc.).
Figure~\ref{fig:category} illustrates the IC violation rate and high-risk
capability prevalence per category.

\begin{figure}[t]
\centering
\includegraphics[width=\columnwidth]{figures/fig_category_ic_breakdown.pdf}
\caption{IC violation rates by skill category.
Left: stacked IC-type distribution.
Right: mean CR vs.\ high-risk capability prevalence (bubble size $\propto$ skill count).
Finance, Security, and AI/Automation occupy the lowest CR, highest risk quadrant.}
\label{fig:category}
\end{figure}

\paragraph{Key findings.}
Three categories form a \emph{high-risk, low-disclosure} cluster
with risk scores $>0.30$: \textbf{Finance}, \textbf{Security}, and \textbf{AI/Automation}.
These categories combine elevated prevalence of dangerous capabilities
(CRED\_READ, NET\_WRITE, AUTONOMOUS) with the lowest mean CR values ($\leq0.20$),
making them the highest-priority targets for marketplace review.

Conversely, \textbf{Education} exhibits the lowest risk score ($0.127$):
its skills tend toward simple content retrieval (NET\_READ) with minimal
credential or file-system access, and their descriptions more accurately
reflect this simpler profile.

Critically, \textbf{IC-1 (Full Omission) is endemic across all categories}---
even categories commonly perceived as low-risk (Monitoring: 16.1\%, Travel: 20.5\%).
This cross-category universality supports the interpretation that under-disclosure
is a broad \emph{marketplace norm} rather than confined to inherently
dangerous domains.
